% !TEX root = ../main.tex

\section{Experimental Design}
\label{sec:design}

\subsection{Datasets}
\label{sec:datasets}

% !TEX root = ../main.tex
% Rows generated from fastener-experiments distilled logs (splits), DATASETS_USED.md and
% HYPOTHESES.md (classes) and h2h3_budgets.json (budgets); see evidence/ for provenance.
\begin{table}[!htbp]
\caption{Datasets, ordered by the number of features $d$. Train and test give rows; the validation part has as many rows as the test part. Budget: model fits per run}
\label{tab:datasets}
\centering
\footnotesize
\setlength{\tabcolsep}{3.5pt}
\begin{tabular}{@{}>{\raggedright\arraybackslash}p{31mm}lrrrrr@{}}
\toprule
Dataset & Domain & $d$ & Train & Test & Classes & Budget \\
\midrule
PenDigits & handwriting & 16 & 6{,}594 & 2{,}199 & 10 & 2{,}000 \\
ItalyPowerDemand & electricity & 24 & 658 & 219 & 2 & 2{,}000 \\
satimage$^{a}$ & satellite & 36 & 3{,}858 & 1{,}286 & 6 & 6{,}000 \\
Crop & satellite & 46 & 14{,}400 & 4{,}800 & 24 & 4{,}000 \\
ProximalPhalanx\-OutlineCorrect & image outline & 80 & 535 & 178 & 2 & 12{,}000 \\
TwoLeadECG & ECG & 82 & 698 & 232 & 2 & 2{,}000 \\
Libras & gesture & 90 & 210 & 75 & 15 & 12{,}000 \\
ElectricDevices & electricity & 96 & 9{,}983 & 3{,}327 & 7 & 6{,}000 \\
MedicalImages & image & 99 & 683 & 229 & 10 & 12{,}000 \\
ECG5000 & ECG & 140 & 3{,}000 & 1{,}000 & 5 & 6{,}000 \\
Wafer & industrial sensor & 152 & 4{,}300 & 1{,}432 & 2 & 2{,}000 \\
ChlorineConcentration & simulation & 166 & 2{,}585 & 861 & 3 & 4{,}000 \\
Adiac & image outline & 176 & 469 & 156 & 37 & 12{,}000 \\
EOData$^{a}$ & satellite & 182 & 29{,}976 & 10{,}008 & 24 & 6{,}000 / 2{,}000 \\
LSST & astronomy & 216 & 2{,}955 & 985 & 14 & 6{,}000 \\
Strawberry & spectra & 235 & 591 & 196 & 2 & 6{,}000 \\
USPS$^{a}$ & handwriting & 256 & 5{,}578 & 1{,}860 & 10 & 6{,}000 \\
Coffee & spectra & 286 & 34 & 11 & 2 & 12{,}000 \\
Ham & spectra & 431 & 128 & 43 & 2 & 12{,}000 \\
FordA & engine sensor & 500 & 2{,}953 & 984 & 2 & 6{,}000 \\
BeetleFly & image outline & 512 & 24 & 8 & 2 & 12{,}000 \\
Earthquakes & seismic & 512 & 275 & 93 & 2 & 12{,}000 \\
har$^{a}$ & motion sensor & 561 & 6{,}181 & 2{,}059 & 6 & 6{,}000 \\
OliveOil & spectra & 570 & 36 & 12 & 4 & 2{,}000 \\
BasicMotions & motion sensor & 600 & 48 & 16 & 4 & 2{,}000 \\
Epilepsy & motion sensor & 618 & 163 & 56 & 4 & 12{,}000 \\
Computers & electricity & 720 & 300 & 100 & 2 & 12{,}000 \\
UWaveGestureLibraryAll & gesture & 945 & 2{,}686 & 896 & 8 & 6{,}000 \\
StarLightCurves & astronomy & 1{,}024 & 5{,}542 & 1{,}847 & 3 & 2{,}000 \\
COIL20$^{a}$ & object images & 1{,}024 & 880 & 280 & 20 & 6{,}000 \\
NATOPS & motion capture & 1{,}224 & 216 & 72 & 6 & 12{,}000 \\
FingerMovements & EEG & 1{,}400 & 248 & 84 & 2 & 6{,}000 \\
CinCECGTorso & ECG & 1{,}639 & 852 & 284 & 4 & 6{,}000 \\
PhonemeSpectra & speech & 2{,}387 & 4{,}016 & 1{,}326 & 39 & 2{,}000 \\
HandOutlines & image outline & 2{,}709 & 822 & 274 & 2 & 12{,}000 \\
gisette$^{a}$ & handwriting & 5{,}000 & 4{,}200 & 1{,}400 & 2 & 6{,}000 \\
FaultDetectionA & vibration & 5{,}120 & 8{,}184 & 2{,}728 & 3 & 6{,}000 \\
EthanolConcentration & spectra & 5{,}253 & 316 & 104 & 4 & 12{,}000 \\
SelfRegulationSCP1 & EEG & 5{,}376 & 337 & 112 & 2 & 12{,}000 \\
Prostate\_GE$^{a}$ & gene expression & 5{,}966 & 62 & 20 & 2 & 6{,}000 \\
FaceDetection & MEG & 8{,}928 & 5{,}650 & 1{,}882 & 2 & 12{,}000 \\
arcene$^{a}$ & mass spectrometry & 10{,}000 & 120 & 40 & 2 & 12{,}000 \\
Heartbeat & heart sound & 24{,}705 & 245 & 82 & 2 & 12{,}000 \\
EigenWorms & motion & 107{,}904 & 155 & 52 & 5 & 12{,}000 \\
PEMS-SF & traffic & 138{,}672 & 264 & 88 & 7 & 2{,}000 \\
MotorImagery & ECoG & 192{,}000 & 226 & 76 & 2 & 6{,}000 \\
\bottomrule
\end{tabular}

\smallskip
\raggedright
$^{a}$~Dataset of the experiment designed for this comparison. EOData is also in the benchmark; its two budgets are those of the designed experiment and of the benchmark.
\end{table}


We use 46 classification datasets with 16 to 192{,}000 features and 2 to 39 classes (Table~\ref{tab:datasets}). Eight of them come from an experiment designed for this comparison and were chosen to cover feature counts from 36 to 10{,}000. They are EOData, a land-cover dataset built from Sentinel-2 satellite observations~\cite{Sircelj2020EOData} for which FASTENER was developed~\cite{Koprivec2020}; satimage and har from OpenML~\cite{Vanschoren2014OpenML}; and USPS, COIL20, gisette, Prostate\_GE, and arcene from the scikit-feature repository~\cite{Li2017ScikitFeature}. The other 38 datasets come from the UCR and UEA time series archives~\cite{Dau2019UCR,Bagnall2018UEA} and were run, together with EOData, in a broader benchmark of feature-selection methods. They were chosen to cover the range of dataset shapes in the archives: the archive datasets were grouped by the logarithms of their numbers of examples and features with $k$-means ($k = 7$), and three to seven datasets were taken from each group.

In the time series datasets, one row is one series, and its time points are the features. For a series with several channels, each pair of a channel and a time point is a feature. Neighboring time points are strongly correlated, so these datasets contain many redundant features, which is the situation feature selection is meant to address.

Every dataset is split once into stratified training, validation, and test parts in the proportions 60/20/20. The split is the same for all seeds and configurations. For the time series, we merge the training and test parts of the archive before splitting, so published accuracies on these datasets are not comparable with ours. EOData has 480{,}000 rows, of which we use a fixed stratified sample of 50{,}000. The search and the choice of the best subset at each size use only the validation part; the test part is used once, for the reported scores. The split of EOData is random rather than spatial, because the published data contain no parcel identifiers. Neighboring pixels of one field can therefore fall into different parts, and its test scores are likely optimistic.

Each run has a fixed budget of model fits. A subset that was evaluated before is taken from a cache of stored results and does not count against the budget. In the designed experiment, the budget is 6{,}000 fits, and 12{,}000 on arcene. In the benchmark, the budget was set for each dataset before the runs, from a calibration run of the FASTENER method on the validation part with a separate seed. It is the smallest of 1{,}000, 2{,}000, 4{,}000, 6{,}000, and 12{,}000 fits at which the calibration run's AUF was within 0.005 of its value at 12{,}000 fits and no more than 0.005 below that of a Boruta reference, increased by one step. This gives 2{,}000 to 12{,}000 fits. The budgets therefore follow the convergence of a different search method, not of NSGA-II.

For comparisons of test scores at matched subset sizes, we use only datasets with more than 50 test rows. With fewer rows, a single misclassified row changes F1 by more than 0.02, the margin we use for these comparisons. This excludes BeetleFly (8 test rows), Coffee (11), OliveOil (12), BasicMotions (16), Prostate\_GE (20), arcene (40), and Ham (43) from these comparisons. All other analyses use every dataset. We set this rule after a first analysis that included BeetleFly and Coffee; with them, the margin is exceeded under the random forest.

\subsection{Search Configurations}
\label{sec:configurations}

We compare three configurations of NSGA-II~\cite{Deb2002NSGAII}. All use the operators of NSGAII-MIIP~\cite{Li2025} in our own reimplementation, written from the authors' published description and code. They differ only in the initial population and in whether the improvement phase runs:
\begin{itemize}
  \item \textbf{Opposition-based start}: the opposition-based initialization of NSGAII-MIIP (Section~\ref{sec:access}), without the improvement phase.
  \item \textbf{Sparse start}: subset sizes drawn uniformly from 1 to 30, then that many features drawn at random, without the improvement phase.
  \item \textbf{Complete NSGAII-MIIP}: the opposition-based start with the mutual-information improvement phase. This configuration was run only in the benchmark.
\end{itemize}

The remaining settings are shared. The population has 100 subsets. Parents are chosen by binary tournament on front rank and then crowding distance. Each pair of parents is crossed with probability 0.9 by single-point crossover at a position where they differ. Each child is then changed by the balanced mutation of Section~\ref{sec:mutation} at the rate $1/d$. The next population is chosen from parents and children together by nondominated sorting and crowding distance. In the complete NSGAII-MIIP, the improvement phase runs in every generation on the current nondominated subsets. It discretizes each feature into ten equal-width bins, estimates mutual information between features and with the class, groups the features into clusters, and adds, removes, or exchanges features according to their relevance to the class and their redundancy with the selected features~\cite{Li2025}. The subsets it produces compete with the others in the next selection.

Every new subset costs one model fit. All fits count against the budget of Section~\ref{sec:datasets}, including those of the improvement phase. A subset that was evaluated before is taken from the cache and costs nothing, so a run can propose far more subsets than it fits. A run ends when its budget is spent. In the benchmark runs with the random forest and LightGBM, a run also ends after 50 consecutive generations that produce no new subset. Without this rule, a run whose population produces only repeated subsets could continue for hundreds of thousands of generations without spending its budget. The runs with the decision tree have no such rule. Every benchmark run was limited to six hours, and runs that reached the limit are counted as missing. Twelve combinations of dataset and configuration were not run, because their estimated memory or time exceeded the available resources. We report how missing runs affect each result.

Each configuration is run with five seeds, 2020 to 2024, on every dataset. The seed changes the search, not the data split. In the designed experiment, the classifier is a decision tree. In the benchmark, every dataset and configuration is run three times, each time with a different classifier: a decision tree, a random forest with 100 trees, and LightGBM. We treat the three as separate groups of runs. The tree and the forest use the scikit-learn defaults. LightGBM uses its defaults except for the smallest number of rows in a leaf, which is a twentieth of the training rows, between 1 and 20, so that the small datasets can still be split. All three use a fixed random state, so the same subset always receives the same score. The classifier both guides the search and produces the reported scores.

\subsection{Statistical Analysis}
\label{sec:statistics}

Datasets are the units of analysis. Seeds of the same dataset share its data split and are not independent samples of the population of datasets, so for each dataset and configuration we first average over seeds; reach becomes the share of seeds that reach the limit. We analyze each group of runs (the designed experiment and the three benchmark groups) separately.

To compare two configurations, we take their paired differences over the datasets on which both have finished runs. We report the numbers of wins, ties, and losses, the mean difference with a 95\% percentile bootstrap interval over datasets (10{,}000 resamples), a two-sided Wilcoxon signed-rank test, and the matched-pairs rank-biserial correlation as the effect size. Because we run several related tests, we correct their $p$-values with the Holm method~\cite{Holm1979}. It controls the probability of at least one false finding across a family of $m$ tests, like the Bonferroni correction, but is less conservative. The $p$-values are sorted from smallest to largest; the smallest is multiplied by $m$, the next by $m - 1$, and so on, and each corrected value is raised, if necessary, to the largest corrected value before it, so that the order is kept. A result counts as significant if its corrected value is below 0.05. A family is a set of tests that answer one question: the four groups of runs for the main comparison of the two starts, and the three pairs of configurations within a group for the other comparisons. Quality at matched sizes (Section~\ref{sec:reach}) counts as equal if its 95\% interval lies within $\pm 0.02$ of zero.

For the time to reach, we use the share of the fit budget spent when a run's archive first holds a subset with at most ten features. Runs that never reach the limit are censored at the full budget. We show Kaplan--Meier curves of the share of runs that have reached the limit and compare configurations by the restricted mean time to reach, up to the full budget, per dataset. To relate reach to the number of features, we fit a logistic regression of reach on $\log_2 d$ and $\log_2$ of the budget over runs, with bootstrap intervals over datasets, and report the number of features at which half of the runs reach the limit at the median budget.

Runs that did not finish are left out of the main analysis. We repeat every comparison in three ways: with only the datasets on which all five seeds finished, with every missing run counted as reaching nothing and scoring zero, and with only the datasets on which every configuration has runs. To compare the classifiers, we use the 26 datasets on which both starts have runs under all three, with a Friedman test and Holm-corrected pairwise Wilcoxon tests. The five seeds per dataset enter only through these averages, and most datasets come from one archive, so the tests describe these datasets and should not be read as estimates for feature selection in general.

\subsection{Pilot Experiment}
\label{sec:pilot}

After the main analysis, we ran a pilot to examine three questions it raised: whether a start that covers all subset sizes keeps the reach of the sparse start without losing the large subsets, whether the opposition-based start's lack of reach is a matter of budget, and whether a mutation that always changes a subset removes the repeats. We registered its criteria before the runs. The pilot uses four datasets from the range of feature counts in which reach changes (Section~\ref{sec:res-d}): Adiac (176 features), Strawberry (235), Earthquakes (512), and Computers (720), each with at least 90 test rows. Their earlier results had been seen, so the pilot guides the interpretation and does not confirm it. It uses the decision tree, five new seeds (2025 to 2029), the benchmark budgets (6{,}000 fits on Strawberry, 12{,}000 on the others), and the stopping rule of 50 generations without a new subset in every run.

The pilot compares ten configurations of NSGA-II without the improvement phase. Four starts are combined with two mutations: the opposition-based, sparse, uniform, and log-uniform starts of Section~\ref{sec:access}, each with the balanced mutation and with a \emph{forced flip}. The forced flip applies the balanced mutation and then, to a child that it left unchanged, adds or removes one feature with equal probability, so the expected change in size stays zero. The two other configurations give the opposition-based start three and ten times the budget. The 1$\times$ and 3$\times$ runs share their seeds and random draws with the 10$\times$ runs, so the three budgets are points on the same search trajectories, not independent runs. Configurations that differ only in mutation start from the same initial population, which we checked for every pair. Every run records each evaluated subset and, per generation, the numbers of proposals, cache hits, and children left unchanged by mutation.

A comparison holds on a dataset if it holds in at least four of five seeds, paired by seed, and a part of a criterion holds if it holds on at least three of the four datasets. Four datasets are too few for a test over datasets.
